%%% BEGIN UNIT
\wbunit{Comparar y ordenar fracciones}{Vas a comparar fracciones para saber cuál es más grande o más pequeña usando dibujos y rectas.}

\section{Comparar con el mismo denominador}

En la fiesta de Ana hay dos pizzas del mismo tamaño. Las dos están cortadas en 8 partes iguales. Ana come 3 partes de una pizza y Beto come 5 partes de la otra. ¿Quién comió más pizza?

Las dos pizzas son iguales y todas las partes son iguales. Entonces, quien tiene más partes tiene más pizza. Beto comió más.

\begin{center}
\wbfracbar[6cm]{3}{8}\\[1mm]
Ana: $\frac{3}{8}$\\[4mm]
\wbfracbar[6cm]{5}{8}\\[1mm]
Beto: $\frac{5}{8}$
\end{center}

Como las partes son del mismo tamaño, $\frac{5}{8} > \frac{3}{8}$.

\begin{wbidea}[La regla de las partes iguales]
Cuando dos fracciones tienen el \textbf{mismo denominador}, las partes son del mismo tamaño. Solo tienes que comparar los \textbf{numeradores}: la fracción con el numerador más grande es la mayor.
\end{wbidea}

\begin{wbexample}[Compara $\frac{3}{5}$ y $\frac{4}{5}$]
\textbf{Paso 1.} Miro los denominadores: son 5 y 5. Son iguales, así que las partes son del mismo tamaño.\\[1mm]
\textbf{Paso 2.} Comparo los numeradores: $4 > 3$.\\[1mm]
\textbf{Paso 3.} Escribo la respuesta: $\frac{4}{5} > \frac{3}{5}$.
\end{wbexample}

\begin{wbexample}[Compara $\frac{5}{8}$ y $\frac{3}{8}$ con dibujos]
Dibujo dos barras del mismo largo y las corto en 8 partes iguales.

\begin{center}
\wbfracbar[6cm]{5}{8}\\[1mm]
$\frac{5}{8}$\\[3mm]
\wbfracbar[6cm]{3}{8}\\[1mm]
$\frac{3}{8}$
\end{center}

En la primera barra hay 5 partes coloreadas. En la segunda hay 3 partes coloreadas. Las partes son del mismo tamaño, y $5 > 3$. Entonces $\frac{5}{8} > \frac{3}{8}$.
\end{wbexample}

\begin{wbtip}[Un error común]
No cambies el denominador ni sumes nada. Si los denominadores son iguales, la respuesta está en los numeradores.
\end{wbtip}

\begin{wbpractice}
\begin{enumerate}
\item Escribe $>$, $<$ o $=$ en cada espacio.
  \begin{enumerate}
  \item $\frac{2}{7}$ \wbblank[1.5cm] $\frac{5}{7}$
  \item $\frac{6}{9}$ \wbblank[1.5cm] $\frac{4}{9}$
  \item $\frac{3}{10}$ \wbblank[1.5cm] $\frac{8}{10}$
  \end{enumerate}
\wblines{2}

\item Colorea $\frac{4}{6}$ en la primera barra y $\frac{2}{6}$ en la segunda. Después escribe $>$ o $<$ entre las dos fracciones.

\begin{center}
\wbfracbar[6cm]{0}{6}\\[1mm]
$\frac{4}{6}$\\[3mm]
\wbfracbar[6cm]{0}{6}\\[1mm]
$\frac{2}{6}$
\end{center}

$\frac{4}{6}$ \wbblank[1.5cm] $\frac{2}{6}$
\wblines{1}

\item ¿Cuál fracción es mayor, $\frac{7}{10}$ o $\frac{3}{10}$? Explica cómo lo sabes.
\wblines{3}

\item Luis comió $\frac{3}{8}$ de un pastel y Sofía comió $\frac{5}{8}$ del mismo pastel. ¿Quién comió más pastel?
\wblines{3}

\item Ordena estas fracciones de menor a mayor: $\frac{1}{4}$, $\frac{3}{4}$, $\frac{2}{4}$.
\wblines{2}

\item Escribe una fracción con denominador 9 que sea mayor que $\frac{4}{9}$.
\wbfracbox
\wblines{1}

\item Explica con tus palabras por qué $\frac{5}{8}$ es mayor que $\frac{2}{8}$.
\wblines{3}
\end{enumerate}
\end{wbpractice}

\section{Comparar con el mismo numerador}

En la panadería hay dos panes del mismo tamaño. Uno se corta en 4 partes iguales y el otro se corta en 8 partes iguales. Tomás toma 1 parte de cada pan. ¿En cuál pan su parte es más grande?

\begin{center}
\wbfracbar[6cm]{1}{4}\\[1mm]
$\frac{1}{4}$\\[3mm]
\wbfracbar[6cm]{1}{8}\\[1mm]
$\frac{1}{8}$
\end{center}

Cuando el pan se corta en 4 partes, cada parte es más grande. Cuando se corta en 8 partes, cada parte es más chica. Entonces $\frac{1}{4} > \frac{1}{8}$.

\begin{wbidea}[La regla de las partes]
Cuando dos fracciones tienen el \textbf{mismo numerador}, la fracción con el \textbf{denominador más pequeño} es la mayor, porque sus partes son más grandes.
\end{wbidea}

\begin{wbexample}[Compara $\frac{1}{4}$ y $\frac{1}{8}$]
\textbf{Paso 1.} Los numeradores son iguales: los dos son 1.\\[1mm]
\textbf{Paso 2.} Miro los denominadores: 4 y 8. Cortar en 4 partes hace partes más grandes que cortar en 8 partes.\\[1mm]
\textbf{Paso 3.} Entonces $\frac{1}{4} > \frac{1}{8}$.
\end{wbexample}

\begin{wbexample}[Compara $\frac{3}{4}$ y $\frac{3}{8}$]
\textbf{Paso 1.} Los numeradores son iguales: los dos son 3.\\[1mm]
\textbf{Paso 2.} Los cuartos son más grandes que los octavos.\\[1mm]
\textbf{Paso 3.} Tres cuartos son más que tres octavos: $\frac{3}{4} > \frac{3}{8}$.

\begin{center}
\wbfracbar[6cm]{3}{4}\\[1mm]
$\frac{3}{4}$\\[3mm]
\wbfracbar[6cm]{3}{8}\\[1mm]
$\frac{3}{8}$
\end{center}
\end{wbexample}

\begin{wbtip}[Cuidado con el 8]
$\frac{1}{8}$ \textbf{no} es mayor que $\frac{1}{4}$, aunque 8 sea mayor que 4. Más partes iguales significa partes más pequeñas.
\end{wbtip}

\begin{wbpractice}
\begin{enumerate}
\item Escribe $>$, $<$ o $=$ en cada espacio.
  \begin{enumerate}
  \item $\frac{1}{3}$ \wbblank[1.5cm] $\frac{1}{6}$
  \item $\frac{2}{5}$ \wbblank[1.5cm] $\frac{2}{10}$
  \item $\frac{1}{4}$ \wbblank[1.5cm] $\frac{1}{4}$
  \end{enumerate}
\wblines{2}

\item Colorea $\frac{1}{2}$ en la primera barra y $\frac{1}{5}$ en la segunda. Después escribe $>$ o $<$ entre las dos fracciones.

\begin{center}
\wbfracbar[6cm]{0}{2}\\[1mm]
$\frac{1}{2}$\\[3mm]
\wbfracbar[6cm]{0}{5}\\[1mm]
$\frac{1}{5}$
\end{center}

$\frac{1}{2}$ \wbblank[1.5cm] $\frac{1}{5}$
\wblines{1}

\item ¿Cuál fracción es mayor, $\frac{3}{4}$ o $\frac{3}{6}$? Explica cómo lo sabes.
\wblines{3}

\item Tomás toma 1 parte de un pan cortado en 4 y 1 parte de un pan cortado en 8. ¿En cuál pan su parte es más grande?
\wblines{3}

\item Completa con $>$ o $<$: $\frac{5}{8}$ \wbblank[1.5cm] $\frac{5}{6}$
\wblines{1}

\item Marca $\frac{1}{4}$ y $\frac{1}{8}$ en las rectas. Después escribe $>$ o $<$ entre las dos fracciones.

\wbnumberline[8cm]{4}{1}{1}
\wbnumberline[8cm]{8}{1}{1}

$\frac{1}{4}$ \wbblank[1.5cm] $\frac{1}{8}$
\wblines{1}

\item Explica con tus palabras por qué $\frac{1}{8}$ es menor que $\frac{1}{4}$.
\wblines{3}
\end{enumerate}
\end{wbpractice}

\section{Ordenar fracciones}

En la clase de deportes hay tres listones del mismo largo. El primero está cortado en 6 partes iguales, el segundo también y el tercero también. Un equipo pintó $\frac{2}{6}$ de un listón, otro pintó $\frac{5}{6}$ y otro pintó $\frac{1}{6}$. ¿Cómo acomodamos las partes pintadas de la más corta a la más larga?

Como los tres listones están cortados en 6 partes iguales, las partes son del mismo tamaño. Solo hay que ordenar los numeradores: 1, 2 y 5.

\begin{center}
\wbfracbar[6cm]{1}{6}\\[1mm]
$\frac{1}{6}$\\[3mm]
\wbfracbar[6cm]{2}{6}\\[1mm]
$\frac{2}{6}$\\[3mm]
\wbfracbar[6cm]{5}{6}\\[1mm]
$\frac{5}{6}$
\end{center}

De menor a mayor: $\frac{1}{6} < \frac{2}{6} < \frac{5}{6}$.

\begin{wbidea}[Cómo ordenar fracciones]
\begin{itemize}
\item Si los \textbf{denominadores son iguales}, ordena por los numeradores.
\item Si los \textbf{numeradores son iguales}, la fracción con el denominador más grande es la más pequeña.
\end{itemize}
\end{wbidea}

\begin{wbexample}[Ordena $\frac{2}{6}$, $\frac{5}{6}$ y $\frac{1}{6}$ de menor a mayor]
\textbf{Paso 1.} Los denominadores son iguales (6, 6 y 6).\\[1mm]
\textbf{Paso 2.} Ordeno los numeradores de menor a mayor: 1, 2, 5.\\[1mm]
\textbf{Paso 3.} Escribo las fracciones en ese orden: $\frac{1}{6} < \frac{2}{6} < \frac{5}{6}$.
\end{wbexample}

\begin{wbexample}[Ordena $\frac{1}{2}$, $\frac{1}{5}$ y $\frac{1}{3}$ de menor a mayor]
\textbf{Paso 1.} Los numeradores son iguales (1, 1 y 1).\\[1mm]
\textbf{Paso 2.} El denominador más grande da las partes más pequeñas. Los denominadores son 2, 5 y 3. El más grande es 5, después 3 y después 2.\\[1mm]
\textbf{Paso 3.} Entonces $\frac{1}{5} < \frac{1}{3} < \frac{1}{2}$.

\begin{center}
\wbfracbar[6cm]{1}{5}\\[1mm]
$\frac{1}{5}$\\[3mm]
\wbfracbar[6cm]{1}{3}\\[1mm]
$\frac{1}{3}$\\[3mm]
\wbfracbar[6cm]{1}{2}\\[1mm]
$\frac{1}{2}$
\end{center}
\end{wbexample}

\begin{wbtip}[Escribe en orden]
Ve escribiendo las fracciones en el mismo orden en que las comparas. Así no se te pierda ninguna.
\end{wbtip}

\begin{wbpractice}
\begin{enumerate}
\item Ordena de menor a mayor: $\frac{3}{8}$, $\frac{1}{8}$, $\frac{6}{8}$.
\wblines{2}

\item Ordena de menor a mayor: $\frac{1}{2}$, $\frac{1}{5}$, $\frac{1}{3}$.
\wblines{2}

\item Ordena de mayor a menor: $\frac{4}{5}$, $\frac{4}{8}$, $\frac{4}{10}$.
\wblines{2}

\item Marca $\frac{2}{6}$, $\frac{5}{6}$ y $\frac{1}{6}$ en la recta. Después escríbelas de menor a mayor.

\wbnumberline[9cm]{6}{1}{1,2,5}
\wblines{2}

\item Tres amigos compartieron barras de chocolate del mismo tamaño. Ana comió $\frac{2}{6}$, Beto comió $\frac{5}{6}$ y Carla comió $\frac{3}{6}$. Ordénalos de menor a mayor según lo que comieron.
\wblines{3}

\item ¿Cuál lista está ordenada de menor a mayor? Encierra la correcta y explica por qué las otras dos no lo están.
  \begin{enumerate}
  \item $\frac{1}{2}$, $\frac{1}{3}$, $\frac{1}{4}$
  \item $\frac{1}{4}$, $\frac{1}{3}$, $\frac{1}{2}$
  \item $\frac{1}{3}$, $\frac{1}{2}$, $\frac{1}{4}$
  \end{enumerate}
\wblines{3}

\item Escribe tres fracciones con denominador 10, ordenadas de menor a mayor.
\wblines{2}
\end{enumerate}
\end{wbpractice}

\begin{wbreview}
\begin{enumerate}
\item Escribe $>$ o $<$: $\frac{3}{7}$ \wbblank[1.5cm] $\frac{6}{7}$
\wblines{1}

\item Escribe $>$ o $<$: $\frac{1}{5}$ \wbblank[1.5cm] $\frac{1}{8}$
\wblines{1}

\item Escribe $>$ o $<$: $\frac{2}{3}$ \wbblank[1.5cm] $\frac{2}{9}$
\wblines{1}

\item Ordena de menor a mayor: $\frac{5}{8}$, $\frac{2}{8}$, $\frac{7}{8}$.
\wblines{2}

\item Ordena de menor a mayor: $\frac{1}{2}$, $\frac{1}{6}$, $\frac{1}{4}$.
\wblines{2}

\item En la tienda venden dos gelatinas del mismo tamaño. Una está cortada en 6 partes iguales y la otra en 4 partes iguales. Rosa compró 3 partes de la gelatina cortada en 6 y Diego compró 3 partes de la gelatina cortada en 4. ¿Quién compró más gelatina? Explica.
\wblines{3}

\item Colorea las barras para mostrar que $\frac{3}{4} > \frac{3}{8}$.

\begin{center}
\wbfracbar[6cm]{0}{4}\\[1mm]
$\frac{3}{4}$\\[3mm]
\wbfracbar[6cm]{0}{8}\\[1mm]
$\frac{3}{8}$
\end{center}
\wblines{1}

\item Explica con tus palabras cómo comparas dos fracciones que tienen el mismo numerador.
\wblines{3}
\end{enumerate}
\end{wbreview}
%%% END UNIT
%%% BEGIN ANSWERS
\wbanswers{\wbLPractice\ 3.1}
\begin{enumerate}
\item a) $\frac{2}{7} < \frac{5}{7}$ \quad b) $\frac{6}{9} > \frac{4}{9}$ \quad c) $\frac{3}{10} < \frac{8}{10}$ (mismo denominador: se comparan los numeradores).
\item Se colorean 4 partes de 6 y 2 partes de 6; $\frac{4}{6} > \frac{2}{6}$.
\item $\frac{7}{10}$ es mayor, porque las dos fracciones tienen décimos del mismo tamaño y $7 > 3$.
\item Sofía comió más: $\frac{5}{8} > \frac{3}{8}$.
\item $\frac{1}{4} < \frac{2}{4} < \frac{3}{4}$.
\item Cualquier fracción con denominador 9 y numerador mayor que 4: $\frac{5}{9}$, $\frac{6}{9}$, $\frac{7}{9}$, $\frac{8}{9}$ o $\frac{9}{9}$.
\item Porque las dos tienen partes del mismo tamaño (octavos) y 5 partes son más que 2 partes.
\end{enumerate}

\wbanswers{\wbLPractice\ 3.2}
\begin{enumerate}
\item a) $\frac{1}{3} > \frac{1}{6}$ \quad b) $\frac{2}{5} > \frac{2}{10}$ \quad c) $\frac{1}{4} = \frac{1}{4}$
\item Se colorea 1 parte de 2 y 1 parte de 5; $\frac{1}{2} > \frac{1}{5}$.
\item $\frac{3}{4}$ es mayor, porque los cuartos son más grandes que los sextos.
\item En el pan cortado en 4 partes, porque $\frac{1}{4} > \frac{1}{8}$.
\item $\frac{5}{8} < \frac{5}{6}$.
\item Se marca $\frac{1}{4}$ y $\frac{1}{8}$ en las rectas; $\frac{1}{4} > \frac{1}{8}$.
\item Porque al cortar en 8 partes iguales cada parte es más pequeña que al cortar en 4 partes iguales.
\end{enumerate}

\wbanswers{\wbLPractice\ 3.3}
\begin{enumerate}
\item $\frac{1}{8} < \frac{3}{8} < \frac{6}{8}$
\item $\frac{1}{5} < \frac{1}{3} < \frac{1}{2}$
\item $\frac{4}{5} > \frac{4}{8} > \frac{4}{10}$
\item En la recta se marcan $\frac{1}{6}$, $\frac{2}{6}$ y $\frac{5}{6}$; $\frac{1}{6} < \frac{2}{6} < \frac{5}{6}$.
\item Ana ($\frac{2}{6}$), Carla ($\frac{3}{6}$), Beto ($\frac{5}{6}$).
\item La correcta es b) $\frac{1}{4} < \frac{1}{3} < \frac{1}{2}$. En a) está al revés: $\frac{1}{2} > \frac{1}{3} > \frac{1}{4}$. En c) no está ordenada porque $\frac{1}{2} > \frac{1}{4}$; con el mismo numerador, entre más grande el denominador, más pequeña es la fracción.
\item Cualquier terna con denominador 10 en orden creciente, por ejemplo $\frac{2}{10} < \frac{5}{10} < \frac{9}{10}$.
\end{enumerate}

\wbanswers{\wbLReview}
\begin{enumerate}
\item $\frac{3}{7} < \frac{6}{7}$
\item $\frac{1}{5} > \frac{1}{8}$
\item $\frac{2}{3} > \frac{2}{9}$
\item $\frac{2}{8} < \frac{5}{8} < \frac{7}{8}$
\item $\frac{1}{6} < \frac{1}{4} < \frac{1}{2}$
\item Diego compró más, porque $\frac{3}{4} > \frac{3}{6}$: los cuartos son más grandes que los sextos.
\item Se colorean 3 partes de 4 y 3 partes de 8; la barra de $\frac{3}{4}$ queda con más sombreado, así que $\frac{3}{4} > \frac{3}{8}$.
\item Respuesta abierta: se comparan los denominadores; la fracción con el denominador más pequeño es la mayor, porque sus partes son más grandes.
\end{enumerate}
%%% END ANSWERS